\section{Funciones Principales y Caracteriisticas de Polars}

\subsection{¿que es Polars?}

Polars es una biblioteca moderna de manipulacion y analisis de datos en Python, diseñada para ofrecer alto rendimiento y eficiencia en el procesamiento de grandes volumenes de informacion. Surge como una alternativa a pandas, proporcionando mayor velocidad y mejor uso de memoria.

Esta desarrollada en el lenguaje Rust, lo que le permite aprovechar la ejecucion paralela y una gestion de memoria optimizada.

\subsection{Características Principales}

\subsubsection{Alto Rendimiento}

\begin{itemize}
    \item Procesamiento más rápido en comparación con bibliotecas tradicionales.
    \item Ejecución paralela automática.
    \item Optimización interna de consultas.
\end{itemize}

\subsubsection{Eficiencia en Memoria}

\begin{itemize}
    \item Basado en Apache Arrow.
    \item Reduce el consumo de memoria RAM.
    \item Permite trabajar con datasets de gran tamaño.
\end{itemize}

\subsubsection{Ejecución Lazy}

Polars permite trabajar en dos modos:

\begin{itemize}
    \item \textbf{Eager:} Ejecuta las operaciones inmediatamente.
    \item \textbf{Lazy:} Construye un plan de ejecución optimizado antes de ejecutarlo.
\end{itemize}

El modo Lazy mejora significativamente el rendimiento en análisis complejos.

\subsubsection{Compatibilidad}

\begin{itemize}
    \item Lectura de archivos CSV, Parquet, JSON y Excel.
    \item Integración con herramientas de visualizacion.
    \item Compatible con entornos como Jupyter Notebook.
\end{itemize}

\subsection{Funciones Principales}

\subsubsection{Lectura de Datos}

\begin{itemize}
    \item \texttt{pl.read\_csv()}
    \item \texttt{pl.read\_parquet()}
    \item \texttt{pl.read\_excel()}
\end{itemize}

\subsubsection{Exploracion de Datos}

\begin{itemize}
    \item \texttt{df.head()}
    \item \texttt{df.tail()}
    \item \texttt{df.describe()}
    \item \texttt{df.schema}
    \item \texttt{df.shape}
\end{itemize}

\subsubsection{Limpieza de Datos}

\begin{itemize}
    \item \texttt{df.drop\_nulls()}
    \item \texttt{df.fill\_null()}
    \item \texttt{df.drop()}
    \item \texttt{df.rename()}
    \item \texttt{df.unique()}
\end{itemize}

\subsubsection{Transformacion y Analisis}

\begin{itemize}
    \item \texttt{df.select()}
    \item \texttt{df.filter()}
    \item \texttt{df.groupby()}
    \item \texttt{df.sort()}
    \item \texttt{df.with\_columns()}
    \item \texttt{pl.col().mean()}
    \item \texttt{pl.col().sum()}
    \item \texttt{pl.col().median()}
    \item \texttt{pl.col().std()}
\end{itemize}

\subsection{Ventajas en Proyectos de Business Intelligence}

\begin{itemize}
    \item Permite analizar grandes volumenes de datos de forma eficiente.
    \item Mejora el rendimiento en comparacion con herramientas tradicionales.
    \item Facilita la creacion de metricas e indicadores (KPIs).
    \item Optimiza consultas mediante ejecucion diferida.
\end{itemize}